% Copyright (c) 2026 Geovana Neves. Licensed under CC BY 4.0 (https://creativecommons.org/licenses/by/4.0/). See LICENSE-AND-AUTHORSHIP.md.
%
% 02-gui-to-pyfs/text/06-the-run.tex
%
% Transcribed from the tables of docs/gui-to-pyfs.md, one row per GUI step:
% the key that takes it (links and solver commands dropped) and the builds
% its commands are verified on. When the page changes a row, change the same
% row here, and the counts of 01-reading.tex with it.

\section{The run}

\begin{frame}{The run: what a key takes (1/2)}
\relax
{\ftstablesize\renewcommand{\arraystretch}{1.05}
\begin{tabularx}{\linewidth}{@{}JKZ@{}}
\guihead
Start the solver & \code{pyfs-matrix plan}, then \code{pyfs-matrix run} & 26.101 to 26.124 \\
Close FlightStream when the script ends & every run type & 26.100 to 26.124 \\
Choose the build & the \code{FS\_BUILD} column & none \\
Run on this machine & \code{pyfs-matrix run}; on Linux with an HPC profile, \code{pyfs-matrix run --local} & none \\
Submit to a cluster & \code{pyfs-matrix run} on Linux with an HPC profile under \code{inputs/hpc/}, then \code{pyfs-matrix collect}; row keys \code{NCPUS} and \code{WALLTIME} & none \\
Stop an unsteady run before its wall clock does, so it can be continued & row key \code{WALLTIME}, with its unit; setup key \code{walltime\_margin\_s} & none \\
Continue an unsteady run where it stopped & row key \code{RESTART}: \code{\{FINISH\_PENDING\}}, \code{\{ADDITIONAL\_ITERS=n\}} or \code{\{ADDITIONAL\_REVS=n\}} & 26.100 to 26.124 \\
\bottomrule
\end{tabularx}}
\guisource{The run}
\end{frame}

\begin{frame}{The run: what a key takes (2/2)}
\relax
{\ftstablesize\renewcommand{\arraystretch}{1.05}
\begin{tabularx}{\linewidth}{@{}JKZ@{}}
\guihead
Limit one point's wall clock on this machine & setup key \code{timeout\_s}, which pyfs enforces around the solver process & none \\
Run only the points not run yet & \code{pyfs-matrix run --resume} & none \\
Redo a point already run & \code{pyfs-matrix run --force-rerun \textless{}point\textgreater{}} & none \\
Skip or refuse a family the mesh does not carry & \code{--ignore-missing-families} of \code{plan} and \code{run}, which reaches each case as row key \code{IGNORE\_MISSING\_FAMILIES} & none \\
\bottomrule
\end{tabularx}}
\par\vspace{1mm}
\begin{takeaway}
Every GUI step of the run has a key: the run is the one stage with no \emph{not yet}.
\end{takeaway}
\guisource{The run}
\end{frame}
